% GENERATED FILE. Edit canonical lessons, then run npm run generate:books.
% canonical-source-hash: fnv1a64:ce4ec1f0e79321ee
% canonical-lessons: TA-C189-pacai, TA-C189-melukuvartti, TA-C189-inaiyatalam, TA-C189-ceyali, TA-C189-tirai

\chapter{Glue, Candle, Website, App, Screen}
\label{ch:ta-glue-candle-website-app-screen}
\hlchaptermodality{189}

\begin{chapteropening}
\textbf{By the end of this chapter:} \emph{I can say glue, a candle, a website, an app and a screen.}

Complete the last lesson of chapter 189: I can say glue, a candle, a website, an app and a screen.
\end{chapteropening}

\section[\texorpdfstring{\ta{பசை} (pacai)}{பசை (pacai)}]{\ta{பசை} (pacai) — glue}

\label{lesson:TA-C189-pacai}



\begin{quote}
Before the new one: say the Tamil for a neem tree, then the Tamil for a coconut palm.
\end{quote}

\subsection*{You'll want to know: \ta{பசை}}
\textbf{\ta{பசை}} — \emph{pacai} — \textquotedblleft{}glue\textquotedblright{}.

\begin{cousinweb}
\textbf{\ta{பசை}} (\emph{pacai}) is native Dravidian. DEDR 3827 lists Malayalam \emph{paca}, \textquotedblleft{}gum, glue\textquotedblright{}.
\end{cousinweb}

\begin{grammarlens}[title={What you've built}]
One more word for this chapter.
\end{grammarlens}

\subsection*{Your turn}
\begin{itemize}
\raggedright
  \item \emph{Say it:} \emph{pacai}
  \item \emph{Say it:} \emph{pacai}, once more
  \item \emph{Say it:} say \emph{pacai}
  \item \emph{Recall:} say the Tamil for a neem tree, then the Tamil for a coconut palm, then say \emph{pacai} again
\end{itemize}

\begin{tcolorbox}[breakable,colback=teal!4,colframe=teal!35!black,title={Before you move on}]
What does \ta{பசை} mean? (\textquotedblleft{}Glue\textquotedblright{}.) Say it once more.
\end{tcolorbox}
\section[\texorpdfstring{\ta{மெழுகுவர்த்தி} (meḻukuvartti)}{மெழுகுவர்த்தி (meḻukuvartti)}]{\ta{மெழுகுவர்த்தி} (meḻukuvartti) — a candle}

\label{lesson:TA-C189-melukuvartti}



\begin{quote}
Before the new one: say the Tamil for a coconut palm, then the Tamil for glue.
\end{quote}

\subsection*{You'll want to know: \ta{மெழுகுவர்த்தி}}
\textbf{\ta{மெழுகுவர்த்தி}} — \emph{meḻukuvartti} — \textquotedblleft{}a candle\textquotedblright{}.

\begin{cousinweb}
\textbf{\ta{மெழுகுவர்த்தி}} (\emph{meḻukuvartti}) is \emph{meḻuku}, \textquotedblleft{}wax\textquotedblright{}, which DEDR 5082 files with words for smearing and glossing, and \emph{vartti}, \textquotedblleft{}wick\textquotedblright{}, from Sanskrit \emph{varti} (CDIAL 11359): \textquotedblleft{}the wax wick\textquotedblright{}.
\end{cousinweb}

\begin{grammarlens}[title={What you've built}]
One more word for this chapter.
\end{grammarlens}

\subsection*{Your turn}
\begin{itemize}
\raggedright
  \item \emph{Say it:} \emph{meḻukuvartti}
  \item \emph{Say it:} \emph{meḻukuvartti}, once more
  \item \emph{Say it:} say \emph{meḻukuvartti}
  \item \emph{Recall:} say the Tamil for a coconut palm, then the Tamil for glue, then say \emph{meḻukuvartti} again
\end{itemize}

\begin{tcolorbox}[breakable,colback=teal!4,colframe=teal!35!black,title={Before you move on}]
What does \ta{மெழுகுவர்த்தி} mean? (\textquotedblleft{}A candle\textquotedblright{}.) Say it once more.
\end{tcolorbox}
\section[\texorpdfstring{\ta{இணையதளம்} (iṇaiyataḷam)}{இணையதளம் (iṇaiyataḷam)}]{\ta{இணையதளம்} (iṇaiyataḷam) — a website}

\label{lesson:TA-C189-inaiyatalam}



\begin{quote}
Before the new one: say the Tamil for glue, then the Tamil for a candle.
\end{quote}

\subsection*{You'll want to know: \ta{இணையதளம்}}
\textbf{\ta{இணையதளம்}} — \emph{iṇaiyataḷam} — \textquotedblleft{}a website\textquotedblright{}.

\begin{cousinweb}
\textbf{\ta{இணையதளம்}} (\emph{iṇaiyataḷam}) is a modern coinage: \emph{iṇaiyam}, \textquotedblleft{}the internet\textquotedblright{}, from native \emph{iṇai}, \textquotedblleft{}to join, to link\textquotedblright{} (DEDR 457), and \emph{taḷam}, \textquotedblleft{}surface, platform\textquotedblright{}, from Sanskrit \emph{tala}, \textquotedblleft{}surface\textquotedblright{} (CDIAL 5731).
\end{cousinweb}

\begin{grammarlens}[title={What you've built}]
One more word for this chapter.
\end{grammarlens}

\subsection*{Your turn}
\begin{itemize}
\raggedright
  \item \emph{Say it:} \emph{iṇaiyataḷam}
  \item \emph{Say it:} \emph{iṇaiyataḷam}, once more
  \item \emph{Say it:} say \emph{iṇaiyataḷam}
  \item \emph{Recall:} say the Tamil for glue, then the Tamil for a candle, then say \emph{iṇaiyataḷam} again
\end{itemize}

\begin{tcolorbox}[breakable,colback=teal!4,colframe=teal!35!black,title={Before you move on}]
What does \ta{இணையதளம்} mean? (\textquotedblleft{}A website\textquotedblright{}.) Say it once more.
\end{tcolorbox}
\section[\texorpdfstring{\ta{செயலி} (ceyali)}{செயலி (ceyali)}]{\ta{செயலி} (ceyali) — an app}

\label{lesson:TA-C189-ceyali}



\begin{quote}
Before the new one: say the Tamil for a candle, then the Tamil for a website.
\end{quote}

\subsection*{You'll want to know: \ta{செயலி}}
\textbf{\ta{செயலி}} — \emph{ceyali} — \textquotedblleft{}an app\textquotedblright{}.

\begin{cousinweb}
\textbf{\ta{செயலி}} (\emph{ceyali}) is native Dravidian (DEDR 1957). It is a modern coinage on \emph{ceyal}, \textquotedblleft{}action\textquotedblright{}, from \emph{cey}, \textquotedblleft{}to do\textquotedblright{}, with \emph{-i}: \textquotedblleft{}the thing that does\textquotedblright{}.
\end{cousinweb}

\begin{grammarlens}[title={What you've built}]
One more word for this chapter.
\end{grammarlens}

\subsection*{Your turn}
\begin{itemize}
\raggedright
  \item \emph{Say it:} \emph{ceyali}
  \item \emph{Say it:} \emph{ceyali}, once more
  \item \emph{Say it:} say \emph{ceyali}
  \item \emph{Recall:} say the Tamil for a candle, then the Tamil for a website, then say \emph{ceyali} again
\end{itemize}

\begin{tcolorbox}[breakable,colback=teal!4,colframe=teal!35!black,title={Before you move on}]
What does \ta{செயலி} mean? (\textquotedblleft{}An app\textquotedblright{}.) Say it once more.
\end{tcolorbox}
\section[\texorpdfstring{\ta{திரை} (tirai)}{திரை (tirai)}]{\ta{திரை} (tirai) — a screen}

\label{lesson:TA-C189-tirai}



\begin{quote}
Before the new one: say the Tamil for a website, then the Tamil for an app.
\end{quote}

\subsection*{You'll want to know: \ta{திரை}}
\textbf{\ta{திரை}} — \emph{tirai} — \textquotedblleft{}a screen\textquotedblright{}.

\begin{cousinweb}
\textbf{\ta{திரை}} (\emph{tirai}) is native Dravidian. DEDR 3244 lists Malayalam \emph{tira}, \textquotedblleft{}curtain\textquotedblright{} and Telugu \emph{tera}, \textquotedblleft{}screen, curtain\textquotedblright{}.
\end{cousinweb}

\begin{grammarlens}[title={What you've built}]
One more word for this chapter.
\end{grammarlens}

\subsection*{Your turn}
\begin{itemize}
\raggedright
  \item \emph{Say it:} \emph{tirai}
  \item \emph{Say it:} \emph{tirai}, once more
  \item \emph{Say it:} say \emph{tirai}
  \item \emph{Recall:} say the Tamil for a website, then the Tamil for an app, then say \emph{tirai} again
\end{itemize}

\begin{tcolorbox}[breakable,colback=teal!4,colframe=teal!35!black,title={Before you move on}]
What does \ta{திரை} mean? (\textquotedblleft{}A screen\textquotedblright{}.) Say it once more.
\end{tcolorbox}
